% Figure: atmospheric absorption coefficient versus frequency for two humidities.
% Approximate power-law growth, ~ f^2 at low f, transitioning to relaxation peaks.
\begin{figure}[htbp]
\centering
\begin{tikzpicture}
\begin{axis}[
  width=12cm, height=7.5cm,
  xlabel={frequency (Hz)}, ylabel={absorption (dB/km)},
  xmode=log, ymode=log,
  xmin=50, xmax=20000, ymin=0.1, ymax=2000,
  axis lines=left,
  legend pos=north west,
  legend cell align=left,
  every axis x label/.style={at={(current axis.right of origin)}, anchor=west},
  every axis y label/.style={at={(current axis.above origin)}, anchor=south},
  grid=major, grid style={enggray!20},
  tick label style={font=\scriptsize},
  label style={font=\small},
  legend style={font=\scriptsize},
]
% representative dry-air absorption (10% RH), read from data file
\addplot[engred, very thick, mark=*, mark size=1pt] coordinates {
  (63,0.3) (125,0.9) (250,2.4) (500,5.0) (1000,9.0) (2000,21)
  (4000,55) (8000,170) (16000,560)};
\addlegendentry{dry air (10\% RH)}
% representative humid-air absorption (70% RH)
\addplot[engblue, very thick, mark=square*, mark size=1pt] coordinates {
  (63,0.1) (125,0.4) (250,1.1) (500,2.6) (1000,4.5) (2000,8.5)
  (4000,18) (8000,42) (16000,120)};
\addlegendentry{humid air (70\% RH)}
\end{axis}
\end{tikzpicture}
\caption[Atmospheric sound absorption versus frequency]{Representative
atmospheric absorption coefficient in $\si{\decibel\per\kilo\meter}$
versus frequency, at two relative humidities and standard temperature
and pressure. Absorption rises steeply with frequency: low-frequency
sound travels for kilometres while high-frequency sound is gone within
a few hundred metres. Humidity shifts the molecular-relaxation
absorption of nitrogen and oxygen, so the curves cross and reorder
across the band rather than scaling uniformly. This frequency selectivity
is why distant thunder rumbles (the high-frequency crack has been
filtered out by the intervening air) and why outdoor public-address
systems favour the mid band. Values read from
\texttt{data/atmospheric\_absorption.csv}; the precise numbers depend on
temperature, humidity, and pressure through the ISO~9613-1 model.}
\label{fig:vol03ch07:atmospheric-attenuation}
\end{figure}
